\documentclass[11pt,a4paper]{article}
\usepackage[T1]{fontenc}
\usepackage[utf8]{inputenc}
\usepackage{lmodern}
\usepackage[a4paper,margin=23mm]{geometry}
\usepackage{booktabs,longtable,array,xurl}
\usepackage[colorlinks=true,urlcolor=blue]{hyperref}
\newcommand{\code}[1]{{\small\nolinkurl{#1}}}
\newcolumntype{P}[1]{>{\raggedright\arraybackslash}p{#1}}
\setlength{\parindent}{0pt}
\setlength{\parskip}{5pt}
\begin{document}
\begin{center}
{\Large\bfseries Client encryption format details}\par
Companion to Specification 2.6.2 (client and server modes), 31 October 2026
\end{center}

This companion preserves the byte-level requirements previously embedded in
Part 4 of Specification 2.2. The main paper retains the security bindings,
validation rules, custody boundary and failure behavior. This file is not
included in the paper. These layouts are part of the versioned format contract,
not evidence of library interoperability or production qualification: RCP-007
still requires the full signed-wire profile, pinned implementations and exact
positive and negative examples before enablement.

\section*{Payload binding (CRY-003)}
The payload AAD is exactly 60 bytes. Its domain identifier, tenant, object,
crypto version, suite and format also bind the key package. Mutable names,
classifications and envelope generation are excluded. AAD is authenticated,
not confidential.

\begin{longtable}{@{}P{0.15\textwidth}P{0.20\textwidth}P{0.55\textwidth}@{}}
\caption{Payload AAD layout}\label{tab:aad}\\
\toprule Bytes & Size & Content\\\midrule\endfirsthead
\toprule Bytes & Size & Content\\\midrule\endhead
0--7 & 8 & ASCII \code{FGAAD001}\\
8--23 & 16 & Tenant UUID\\
24--39 & 16 & Object UUID\\
40--55 & 16 & Crypto-version UUID\\
56--57 & 2 & Suite identifier, unsigned big-endian, from a versioned registry\\
58--59 & 2 & Format version, unsigned big-endian\\
\bottomrule
\end{longtable}

\section*{Key package inside HPKE (KEY-006)}
The version-1 key package is encoded in this order:
\begin{enumerate}
\item Unsigned big-endian 16-bit format value 1.
\item Unsigned big-endian 16-bit binding length 60.
\item Exactly 60 payload-binding bytes from the preceding layout.
\item Unsigned big-endian 32-bit keyset length.
\item Exactly that many bytes of the library-serialized binary keyset.
\end{enumerate}

The receiving client rejects inconsistent lengths, trailing bytes, unknown
versions and keysets above 64 KiB. It validates the suite, embedded payload
binding and single-key streaming-AEAD template before use. Extra keysets or key
entries reject. The library owns its keyset serialization and cryptographic
operations; this layout does not authorize reimplementing them.

The server sees only the outer HPKE envelope. It cannot validate the inner
keyset or prove that a malicious sender encrypted only the declared material.
The main paper's independently authorized publication and delivery checks
remain mandatory.

\section*{Server mode (D-12, SSE-001, SSE-002)}
A server-mode version has no key package or envelope. Its data key is 32
random bytes generated by the \code{keys} role, which never sees a payload;
the storage role encrypts and decrypts with it for one stream. Two byte
layouts bind it.

\textbf{Sealed data key.} The \code{keys} role seals the data key with
AES-256-GCM (Tink, no output prefix) under the tenant's KEK version: 12-byte
nonce, ciphertext, 16-byte tag. The catalog stores the result with the KEK
version. Its associated data:

\begin{longtable}{@{}P{0.15\textwidth}P{0.20\textwidth}P{0.55\textwidth}@{}}
\caption{Sealed data key associated data}\label{tab:sealaad}\\
\toprule Bytes & Size & Content\\\midrule\endfirsthead
\toprule Bytes & Size & Content\\\midrule\endhead
0--14 & 15 & ASCII \code{fga-kek-seal/v1}\\
15 & 1 & Zero\\
16--19 & 4 & KEK version, unsigned big-endian\\
20--35 & 16 & Tenant UUID\\
36--51 & 16 & Object UUID\\
52--67 & 16 & Logical-version UUID\\
68--83 & 16 & Crypto-version UUID\\
84--99 & 16 & Blob UUID\\
\bottomrule
\end{longtable}

\textbf{Payload.} The storage role encrypts the payload with Tink's
streaming AEAD AES-256-GCM-HKDF (HKDF-SHA256, 32-byte derived keys, 1 MiB
segments), whose key value is the data key itself. The version records suite
1 and format 2. The ciphertext of $n$ plaintext bytes is
$40 + n + 16s$ bytes for $s$ segments, where the first segment holds up to
$2^{20} - 56$ plaintext bytes and each further one $2^{20} - 16$. Its
associated data:

\begin{longtable}{@{}P{0.15\textwidth}P{0.20\textwidth}P{0.55\textwidth}@{}}
\caption{Server-mode payload associated data}\label{tab:serveraad}\\
\toprule Bytes & Size & Content\\\midrule\endfirsthead
\toprule Bytes & Size & Content\\\midrule\endhead
0--20 & 21 & ASCII \code{fga-server-payload/v1}\\
21 & 1 & Zero\\
22--37 & 16 & Tenant UUID\\
38--53 & 16 & Object UUID\\
54--69 & 16 & Logical-version UUID\\
70--85 & 16 & Crypto-version UUID\\
86--101 & 16 & Blob UUID\\
\bottomrule
\end{longtable}

Both domains differ from the client payload's \code{FGAAD001}, so a sealed
key or a payload moved to another identity, or read in the other mode, fails
authentication. The bindings cannot keep a party who holds the KEKs, or the
storage role while it serves a read, from reading the content (N-09).

\end{document}
